\section{はじめに}

本書は初めてSCALEモデルを利用するユーザー向けに、SCALE-RMモデルについて解説したユーザーズガイドです。
本書はSCALE version 0.2.6 (SCALE-RM version 4.2.6)に対応した説明になっています。
第\ref{sec:overview}章でSCALEの概要について説明します。第\ref{sec:install}章で必要な環境、
およびインストール方法を説明します。つづいて、第\ref{sec:tuto_ideal}章では理想実験、
第\ref{sec:tuto_real}章では現実大気実験の簡単な例を挙げて実行方法を説明します。
第\ref{sec:install}章から第\ref{sec:tuto_real}章まではひとつながりの
チュートリアルになっており、前章の経験や結果に基づいて説明する部分がありますので注意してください。
第\ref{sec:advance}章ではチュートリアルでは説明しなかった実行方法の詳細、
機能やツールについて説明します。

本書中の不明点やお気づきの点がございましたら、SCALE user's メーリングリスト\\
 \verb|scale-user@riken.jp|までご連絡ください。



\subsection{SCALEの特徴}
SCALEは計算機を用いて気象・気候科学の問題にアプローチする際に、ユーザーが研究を進めやすいように
事前処理から数値モデル計算、解析に至るまですべての過程を網羅する計算ライブラリの開発を
目指したソフトウェアであり。SCALEは下記に挙げるような特徴を持つ。
\begin{itemize}
\item SCALEは、 ``Scalable Computing for Advanced Library and Environment''の略である．
\item SCALEは、フリーソフトウェアである。「BSD-2ライセンス」においてフリーソフトウェア
として提供されており、商用、非商用に関わらず自由な利用・改変が可能なソフトウェアである。
\item SCALEには、SCALE-RMモデル、SCALE-GMといった組み上げ済みの数値モデルが含まれている。
\item SCALEには、次節で説明する様々なコンポーネントが導入されており、必要に応じて取り替えることが可能性である。
\end{itemize}

ライセンスの詳細は、トップディレクトリ直下の``\verb|scale/LICENSE|''のファイルに記述されている。
SCALEの使用前に一読しておくこと。またSCALEのWebページにも説明が記載されているので必要に応じて
参照すること（\url{http://scale.aics.riken.jp/}）。

\subsection{SCALE-RMモデルの構成}
現在，SCALE-RMモデルとして組み込まれているコンポーネントは下記のものである。
詳細なモデル構成や差分化手法については、\cite{scale_2015}、\cite{satoy_2015b}、および\cite{nishizawa_2015}を参照されたい。\\

\noindent{\bf フレームワーク関係}

\begin{itemize}
 \item Cartesian グリッドシステム（Arakawa C-grid）
 \item Cartesian ベースMPI通信
 \item Map-projection \& Map-factors
 \item Nesting システム（1-way：親→子のみ）
   \begin{itemize}
    \item Option: オフライン実行とオンライン実行を選択可能
   \end{itemize}
 \item 複数事例一括実行 システム（バルクジョブシステム）
 \item gtoolベースnetcdf4ファイル I/O
 \item 外部データ読み込みによる初期値・境界値データ作成
   \begin{itemize}
    \item Option: WRF-ARW、NICAM、その他GrADSフォーマットでの入力に対応
   \end{itemize}
\end{itemize}

{\bf 力学コア関係}

\begin{itemize}
 \item 方程式系: ３次元完全圧縮流体方程式
 \item 数値解法: 陽解法と陰解法の両方を実装済み
   \begin{itemize}
    \item Option: HE-VE、HE-VI、HI-VIスキームから選択可能
   \end{itemize}
 \item 空間差分: 4次中央差分
 \item 時間差分: 3段ルンゲクッタスキーム（2次）
 \item 非負保証: FCTスキーム
 \item 数値フィルター: 4次Hyper diffusion
 \item 地形: Terrain-followingスキーム
\end{itemize}

{\bf 物理過程}

\begin{itemize}
 \item 乱流過程: 複数から選択可能
   \begin{itemize}
    \item Smagorinsky-Lilly typeのサブグリッドモデル (\cite{smagorinsky_1963}; \cite{lilly_1962})
    \item MYNN2.5境界層モデル (\cite{my_1982},\cite{nakanishi_2004})
   \end{itemize}
 \item 雲微物理: 複数から選択可能
   \begin{itemize}
    \item Kessler type バルクモデル(\cite{kessler_1969})
    \item 6-class single moment バルクモデル(\cite{tomita_2008})
    \item 6-class double moment バルクモデル(\cite{sn_2014})
    \item ビン法雲モデル(\cite{suzuki_etal_2010})から選択可能
   \end{itemize}
 \item 放射過程: MSTRN (\cite{sekiguchi_2008})
 \item 地表面モデル
  \begin{itemize}
   \item 陸面モデル: バケツモデル（バルク交換係数は \cite{beljaars_1991}; \cite{wilson_2001}）
   \item 海洋モデル:（スラブモデル）
   \item 都市モデル: 単層キャノピーモデル(\cite{kusaka_2001})
  \end{itemize}
\end{itemize}

上記に加えて、SCALE-RMモデル本体のドライバー、現実事例の計算に必要な地形や土地利用データを作成するツール、
理想実験の初期値・境界値、もしくは外部モデルから初期値・境界値を作成するツールが提供されている。

次節でSCALEライブラリの思想とモデルの関係について説明するが、SCALE-RMモデルの実行とは直接関係ないため、
必要なければ読み飛ばしても構わない。


\subsection{SCALEのライブラリとモデル}

\begin{figure}[bh]
\begin{center}
  \includegraphics[width=0.9\hsize]{./figure/library.eps}\\
  \caption{SCALEライブラリのねらい}
  \label{fig:scale}
\end{center}
\end{figure}

SCALEは、理化学研究所 計算科学研究機構（RIKEN AICS）を中心に開発が進められている気象・気候科学向けのライブラリである。
図 \ref{fig:scale}はSCALEのねらいを表したポンチ絵である。この図に示されるように、SCALEプロジェクトは、近い将来もしくは
現在すでに研究の足かせとなっている諸問題を解決する、もしくは解決の手助けをすることを目指している。
SCALEは次世代のスーパーコンピュータから汎用計算機に至るまで広く用いられる事を念頭において開発されており、
気候・気象科学を専門とする科学者と計算機科学を専門とする科学者が共同で開発を行っている（\cite{abe_etal_2013}）。
そのため、スーパーコンピュータ「京」やFujitsu FX10等のスーパーコンピュータ上でも、
ユーザーがチューニングすることなく高い計算効率で実行可能なようになっている。

SCALEライブラリは、ユーザーにスーパーコンピュータ「京」をはじめとする並列計算機に対してチューニングされた
モデルコンポーネントや解析システム、テストセット、そしてこれらに対する知見といった「開発環境」を提供する。

提供された開発環境利用の例として数値モデルの構築が挙げられる。その一例として、SCALEライブラリを用いて構築した数値モデル、
SCALE-RMモデルが本ライブラリセットに含まれている(図 \ref{fig:scale-rm})。
格子系、力学コア、そして物理過程といったモデル構築に必要な基本的コンポーネントはSCALEライブラリが提供する。そのため、
「SCALE-RMモデル」として新たに用意されたのは、これらのコンポーネントを動かすためのモデルドライバーや変数セットである。
このSCALE-RMモデルは研究に利用できるものである、すでに層積雲の遷移過程を調べる研究などの実績がある
（\cite{satoy_2014}、\cite{satoy_2015}）。ユーザーは、SCALEライブラリを用いてSCALE-RMモデルと同様な方法によって
ユーザーは新たなモデル作ることができる。


\begin{figure}[tb]
\begin{center}
  \includegraphics[width=0.9\hsize]{./figure/scale_and_scale-rm.eps}\\
  \caption{SCALEライブラリとSCALE-RMモデルの関係}
  \label{fig:scale-rm}
\end{center}
\end{figure}


\section{表記上の注意}

本書中のシェルコマンド等はbash を想定して記述する。異なる環境下では適宜読み替えて対応すること。
また、本書内では特に断りがない限り、下記の表記法に従うものとする。


コマンドラインのシンボル（\verb|$, #|）は、コマンドの実行を示す。
下記の表記の違いは、プログラムを実行する権限の違いを示している。

\begin{verbatim}
 #        <- root権限で実行するコマンド
 $        <- ユーザ権限で実行するコマンド
\end{verbatim}
権限の一時的な切り替えにはsuコマンドを用いる。
\verb|{User_Name}|は実際のユーザ名に読み替えること。
\begin{verbatim}
 $ su {User_Name}   <- {User_Name}のユーザー名でログイン
 $ exit             <- {User_Name}のユーザー名でログインを修了
 $ su -             <- root権限に変更
 #
\end{verbatim}

コマンドオプションにハイフンを用いると、そのユーザでのログインを行う。
用いない場合、権限のみの変更となる。またユーザ名を省略するとrootでのログインを試す。
ユーザの一時切り替えを終わるには、exitコマンドを用いる。
%各プログラムをインストールするための圧縮ファイルは、/tmpにダウンロードされていると仮定する。
%他のディレクトリにダウンロードしてある場合は、mvコマンド等を用いて/tmpに移動しておくことを勧める。
文章表記のうち、ダブルスラッシュ（//）で始まる行は解説のためのもので、実際に記述する必要はない。

下記のように四角い囲みで区切られた部分は、コマンドラインのメッセージ部分であることを意味する。\\

\noindent {\small {\gt
\fbox{
\begin{tabularx}{140mm}{l}
 -- -- -- -- コマンドラインのメッセージ\\
 -- -- -- -- -- -- -- -- コマンドラインのメッセージ\\
 -- -- -- -- -- -- -- -- -- -- -- -- コマンドラインのメッセージ\\
\end{tabularx}
}}}\\

また、下記のように丸い囲みで区切られた部分は、エディタでファイルを編集する記述内容を表す、
もしくはファイル内の記述を参照している部分である。\\

\noindent {\small {\gt
\ovalbox{
\begin{tabularx}{140mm}{l}
 -- -- -- -- ファイル中の記述\\
 -- -- -- -- -- -- -- -- ファイル中の記述\\
 -- -- -- -- -- -- -- -- -- -- -- -- ファイル中の記述\\
\end{tabularx}
}}}\\


\begin{verbatim}
 $ vi
\end{verbatim}

上記のコマンドは、vi（汎用テキストエディタ）プログラムを実行する。
他に、gedit、emacsなどのテキストエディタがある。
マニュアルではviコマンドの実行指示があるが、各自の使いやすいエディタへ適宜読み替えること。


